\chapter{Что делают нейроны во сне}
\chaptermark{Сон}
\label{sec:sleep}

\section{Сон --- не выключение, а другой режим}

Инженер, увидев выключенный компьютер, скажет: он ничего не делает. Со спящим мозгом так не выходит. Нейроны во сне не молчат: в глубоком сне кора выдаёт импульсы, в быстром сне --- почти столько же, сколько днём. Меняется не то, \emph{работают} ли нейроны, а то, \emph{как} они работают. Сон --- это набор режимов машины, и переключают их те же широковещательные каналы, что и днём.

Для каждого режима можно выписать «состояние регистров» (таблица~\ref{tab:sleepmodes}). Днём \nt{ACET}, \nt{NORA} и \nt{SERO} работают, датчики подключены, мышцы слушаются. В медленном, глубоком сне все три канала затихают, а вход от органов чувств ослаблен~\cite{Hasselmo1999}: выброс ацетилхолина в коре падает, а \nt{NORA}- и \nt{SERO}-нейроны выдают импульсы реже, чем днём~\cite{Marrosu1995,AstonJonesBloom1981,McGintyHarper1976}. В быстром сне картина странная: \nt{ACET} снова поднимается до дневного уровня~\cite{Marrosu1995}, а \nt{NORA} и \nt{SERO} замолкают почти полностью~\cite{AstonJonesBloom1981,McGintyHarper1976,JacobsAzmitia1992}. Мышцы тела в быстром сне выключены: \nt{ACET}-нейроны, управляющие ими, сильно тормозятся через \nt{GABA} и \nt{GLYC}~\cite{BrooksPeever2012} --- тем же запретом из главы~\ref{sec:gaba}, только наложенным на весь выход сразу.

\begin{table}[t]
\centering
\footnotesize
\setlength{\tabcolsep}{4pt}
\renewcommand{\arraystretch}{1.15}
\begin{tabular}{>{\raggedright}p{3.1cm}>{\raggedright}p{2.9cm}>{\raggedright}p{3.2cm}>{\raggedright\arraybackslash}p{3.4cm}}
\toprule
& бодрствование & медленный сон & быстрый сон \\
\midrule
\nt{ACET} (запись, гл.~\ref{sec:acet}) & высокий & низкий & высокий \\
\nt{NORA} (усиление, гл.~\ref{sec:nora}) & средний, вспышки & низкий & почти молчит \\
\nt{SERO} (гл.~\ref{sec:serocomputer}) & средний & низкий & почти молчит \\
вход от датчиков & подключён & ослаблен & ослаблен \\
выход на мышцы & подключён & ослаблен & выключен торможением \\
активность коры & ровная & медленные качели: работа --- тишина & ровная, как днём \\
\bottomrule
\end{tabular}
\caption{Состояние регистров машины в трёх режимах. Все строки измерены; уровни \nt{ACET}, \nt{NORA} и \nt{SERO} --- прямыми записями по стадиям сна~\cite{Marrosu1995,AstonJonesBloom1981,McGintyHarper1976}.}
\label{tab:sleepmodes}
\end{table}

\section{Когда спать: реле с гистерезисом}

Что решает, когда машине переключаться в сон? Хорошо работает простая схема из двух частей~\cite{Borbely1982}. Первая часть --- \emph{давление сна} $S$, которое копится, пока мы бодрствуем, и сбрасывается, пока спим. Это конденсатор, который заряжается днём и разряжается ночью:
\begin{empheq}[box=\keybox]{equation}
\frac{\dd S}{\dd t} \;=\; \frac{1-S}{\tau_r}\ \ \text{(бодрствование)},\qquad
\frac{\dd S}{\dd t} \;=\; -\frac{S}{\tau_d}\ \ \text{(сон)} .
\label{eq:sleeppressure}
\end{empheq}
Вторая часть --- два порога: машина засыпает, когда $S$ доходит до верхнего порога $H(t)$, и просыпается, когда $S$ опускается до нижнего $L(t)$. Оба порога плавно качаются в такт суточному ритму из главы~\ref{sec:chemoutput}. Получилось реле с гистерезисом --- триггер Шмитта из раздела~\ref{sec:bistable}, --- а вместе с конденсатором оно даёт релаксационный генератор, подстраиваемый по фазе суточным ритмом~\eqref{eq:pll}.

\begin{figure}[t]
\centering
\begin{tikzpicture}[>=Stealth,x=0.16cm,y=4.2cm]
\foreach \a/\b in {0/4.3,21.2/29.3,45.4/53.4,69.5/72}{\fill[blue!8] (\a,0) rectangle (\b,1);}
\draw[->,gray!70!black] (0,0) -- (75,0) node[right,font=\small,black] {$t$, ч};
\draw[->,gray!70!black] (0,0) -- (0,1.08) node[left,font=\small,black] {$S$};
\foreach \t in {0,24,48,72}{\draw[gray!60!black] (\t,-0.02) -- (\t,0.02); \node[below,font=\scriptsize] at (\t,0) {\t};}
\draw[thick,densely dashed,red!70!black] plot file {figures/sleep_H.dat};
\draw[thick,densely dashed,green!45!black] plot file {figures/sleep_L.dat};
\draw[very thick,blue!60!black] plot file {figures/sleep_S.dat};
\node[font=\scriptsize,red!70!black,anchor=west] at (73,0.72) {$H$: заснуть};
\node[font=\scriptsize,green!45!black,anchor=west] at (73,0.24) {$L$: проснуться};
\node[font=\scriptsize,blue!60!black] at (25.25,0.93) {сон};
\node[font=\scriptsize,blue!60!black] at (49.4,0.93) {сон};
\end{tikzpicture}
\caption{Давление сна $S$~\eqref{eq:sleeppressure} при $\tau_r=18.2$~ч, $\tau_d=4.2$~ч. Днём $S$ заряжается до верхнего порога $H$, ночью (закрашено) разряжается до нижнего $L$; оба порога качаются вместе с суточным ритмом. Машина сама приходит к режиму около $16$ часов бодрствования и $8$ часов сна.}
\label{fig:twoprocess}
\end{figure}

В нашей модели с обычными значениями $\tau_r=18.2$~ч и $\tau_d=4.2$~ч машина сама выходит на $16.1$ часа бодрствования и $8.0$ часа сна (рис.~\ref{fig:twoprocess}). Модель объясняет и то, что кажется странным: после сорока часов без сна давление доходит до $0.90$, но восстановительный сон в модели длится лишь $8.8$ часа, а не на сутки дольше обычного. Разряд экспоненциальный, и высокий заряд уходит быстро; «долг» возвращается не длительностью, а глубиной сна.

Что физически играет роль $S$? Сильный кандидат --- аденозин, «счётчик нагрузки» из приложения~\ref{sec:othermediators}: его концентрация растёт за время бодрствования и падает во сне~\cite{PorkkaHeiskanen1997,Dunwiddie2001}. Что давление сна --- это именно аденозин и ничего больше, --- гипотеза.

\section{Медленный сон: вся сеть как релаксационный генератор}

В глубоком сне нейроны коры работают по очереди всей группой: реже раза в секунду они вместе включаются на несколько сотен миллисекунд (\emph{состояние работы}) и вместе замолкают (\emph{состояние тишины}): частота этих качелей ниже $1$~Гц, под наркозом чаще всего $0.3$--$0.4$~Гц~\cite{Steriade1993}. Эти медленные качели мы можем собрать из деталей, которые уже есть. Возьмём группу \nt{GLUT}-нейронов с сильной обратной связью на себя --- память из главы~\ref{sec:glutcomputer}, у которой два устойчивых состояния, --- и добавим медленную усталость, как в генераторе ритма главы~\ref{sec:muscles}:
\begin{empheq}[box=\keybox]{equation}
\tau\,\frac{\dd r}{\dd t} = -r + F\bigl(w\,r + I - a\bigr),
\qquad
\tau_a\,\frac{\dd a}{\dd t} = -a + b\,r .
\label{eq:updown}
\end{empheq}
Группа включается, работает и устаёт; усталость $a$ съедает верхнее состояние, и группа падает в тишину; в тишине усталость проходит, нижнее состояние исчезает, и группа снова включается. Получается тот же релаксационный генератор, что и давление сна, только примерно в восемьдесят тысяч раз быстрее.

\begin{figure}[t]
\centering
\begin{tikzpicture}[>=Stealth,x=2.9cm,y=1.0cm]
\begin{scope}[yshift=1.9cm]
\draw[->,gray!70!black] (0,0) -- (4.1,0);
\draw[->,gray!70!black] (0,0) -- (0,1.25) node[left,font=\small,black] {$r$};
\draw[thick,blue!60!black] plot file {figures/updown_sleep.dat};
\node[font=\scriptsize,anchor=west] at (4.15,0.5) {сон: $b=5$};
\end{scope}
\draw[->,gray!70!black] (0,0) -- (4.1,0) node[right,font=\small,black] {$t$, с};
\draw[->,gray!70!black] (0,0) -- (0,1.25) node[left,font=\small,black] {$r$};
\draw[thick,red!70!black] plot file {figures/updown_wake.dat};
\node[font=\scriptsize,anchor=west] at (4.15,0.5) {день: $b=1$};
\foreach \t in {0,1,2,3,4}{\draw[gray!60!black] (\t,-0.05) -- (\t,0.05); \node[below,font=\scriptsize] at (\t,0) {\t};}
\end{tikzpicture}
\caption{Одна и та же группа~\eqref{eq:updown} в двух режимах: $w=5$, $I=2$, $\theta=2.5$, $\tau=10\ms$, $\tau_a=0.3$~с. С сильной усталостью ($b=5$, как в медленном сне) группа качается между работой и тишиной с периодом $1.1$~с (значение модели; в коре период больше секунды, под наркозом --- около $3$~с). Ослабим усталость ($b=1$) --- так действует ацетилхолин --- и та же группа работает ровно.}
\label{fig:updown}
\end{figure}

А что переводит сеть из качелей в ровную дневную работу? Ацетилхолин подавляет в нейронах медленные токи, создающие усталость~\cite{McCormick1992}. В нашей модели при сильной усталости, $b=5$, группа качается с периодом $1.1$~с --- быстрее измеренных качелей, но того же порядка --- и работает около $35\%$ времени, если усреднять по целому числу периодов; при слабой, $b=1$, та же группа работает ровно (рис.~\ref{fig:updown}). Один регистр --- уровень \nt{ACET} --- переключает генератор в усилитель. Поверх медленных качелей в сне идут и более быстрые вспышки ритма $7$--$15$ колебаний в секунду; их порождает петля \nt{GLUT}--\nt{GABA} между корой и промежуточным релейным узлом~\cite{SteriadeMcCormickSejnowski1993}, родственница ритма из главы~\ref{sec:gaba}.

\section{Повторы: перемотка дня}

В главе~\ref{sec:acet} мы видели, что в медленном сне нейроны, работавшие днём вместе, снова срабатывают вместе и в том же порядке. Добавим одну инженерную деталь: повторы идут \emph{с перемоткой}. Последовательность, занимавшая днём секунды, проигрывается во сне за десятые доли секунды~\cite{Nadasdy1999}: в гиппокампе примерно в двадцать раз быстрее~\cite{LeeWilson2002}, в коре --- в шесть-семь раз~\cite{Euston2007}. И проигрывается не когда придётся: короткие вспышки повторов в одной области подстраиваются под качели работы и тишины и под быстрые вспышки ритма в коре. Если искусственно усилить эту согласованность, память улучшается~\cite{Maingret2016} --- это уже причинная связь, а не совпадение.

Зачем машине перематывать день? Инженеру знакома трудность, которую называют \emph{катастрофическим забыванием}: сеть, которая учит новое подряд, одно за другим, портит старое. Спасает \emph{перемешивание}: учить новое вперемешку со старым, небольшими порциями, много раз. Гипотеза о двух системах обучения~\cite{McClelland1995} состоит в том, что быстрая память записывает день целиком, а во сне понемногу, вперемешку и многократно проигрывает его медленной коре. Это та же пара «быстрый буфер --- медленное хранилище», что в главе~\ref{sec:acet}, и та же пара быстрого и медленного учеников, что в главе~\ref{sec:motorlearning}. Повторы и их согласованность измерены; что их назначение --- перемешанное обучение медленной памяти, --- хорошо обоснованная гипотеза.

\section{Перенормировка весов}

Днём правила Хебба из главы~\ref{sec:dofa} в основном усиливают связи. Если только усиливать, веса растут, и вместе с ними растут расходы энергии и падает чувствительность: нейрон, у которого все входы сильные, перестаёт различать их. По гипотезе \emph{синаптического гомеостаза}~\cite{TononiCirelli2014}, сон нужен в том числе для того, чтобы вернуть веса к рабочему уровню: все связи ослабляются в одно и то же число раз, так что их \emph{отношения} --- выученное --- сохраняются. Для инженера это нормировка, как в правиле~\eqref{eq:oja}, только сделанная не на ходу, а в отдельное время.

Прямые измерения этому не противоречат: у мышей площадь контакта на синапсах коры после сна примерно на $18\%$ меньше, чем после бодрствования, уменьшение пропорционально размеру, а самые крупные и устойчивые контакты почти не меняются~\cite{deVivo2017}. Масштабирование весов измерено; что это главная задача сна, --- гипотеза.

\section{Быстрый сон: машина, отключённая от входов и выходов}

В быстром сне кора работает почти как днём, но вход от органов чувств ослаблен, а выход на мышцы выключен торможением. Для инженера это знакомый режим: \emph{испытание на стенде}. Схема работает на полную мощность, но её входы замкнуты на внутренний генератор, а выходы отсоединены от исполнительных механизмов, чтобы ничего не сломать. Сновидения по одной из гипотез --- именно такая внутренняя прогонка модели мира, собранной за день; что именно сеть при этом делает и учится ли, неизвестно. Измерено здесь только то, что записано в таблице~\ref{tab:sleepmodes}: какой канал включён, какой молчит, и что выход заблокирован.

\section{Обслуживание и местный сон}

Долго считалось, что во сне межклеточные щели расширяются и мозг быстрее промывается от продуктов обмена веществ~\cite{Xie2013}. Недавние опыты, измерявшие промывку другим способом, получили обратное: во сне она медленнее~\cite{Miao2024}. Вопрос сейчас открыт, и мы его оставим открытым.

Надёжнее другой факт: сон --- свойство местных сетей, а не всей машины сразу. Если животное долго не давать спать, то отдельные участки его коры, пока животное бодрствует и двигается, ненадолго проваливаются в тишину, как в медленном сне, и в эти моменты животное чаще ошибается~\cite{Vyazovskiy2011}. Каждая сеть устаёт сама и сама переключается; общий режим получается тогда, когда широковещательные каналы переводят все сети сразу.

\begin{keyblock}
\begin{proposition}[Сон как набор режимов]
\label{prop:sleep}
Во сне нейроны не выключаются: широковещательные каналы переводят машину в другие режимы. Когда спать, решает реле с гистерезисом~\eqref{eq:sleeppressure}, подстроенное суточным ритмом. В медленном сне сеть становится релаксационным генератором~\eqref{eq:updown} и проигрывает день с перемоткой; в быстром --- работает, отключённая от входов и выходов. Режимы, повторы и масштабирование весов измерены; их назначение --- перемешанное обучение и перенормировка --- хорошо обоснованные гипотезы.
\end{proposition}
\end{keyblock}

\section{Итог}

\begin{table}[t]
\centering
\footnotesize
\setlength{\tabcolsep}{4pt}
\renewcommand{\arraystretch}{1.15}
\begin{tabular}{>{\raggedright}p{3.2cm}>{\raggedright}p{4.4cm}>{\raggedright\arraybackslash}p{5.2cm}}
\toprule
Что происходит & Как & Что известно \\
\midrule
смена режимов & регистры \nt{ACET}, \nt{NORA}, \nt{SERO} (табл.~\ref{tab:sleepmodes}) & измерено \\
когда спать & конденсатор $S$ и реле с гистерезисом~\eqref{eq:sleeppressure} & модель хорошо описывает опыты; аденозин как $S$ --- гипотеза \\
медленные качели & группа с обратной связью и усталостью~\eqref{eq:updown} & качели измерены; модель --- простейшая схема \\
перемотка дня & повторы в $6$--$20$ раз быстрее, согласованные с качелями & измерено; усиление согласованности улучшает память \\
зачем повторы & перемешанное обучение медленной памяти & хорошо обоснованная гипотеза \\
перенормировка весов & масштабирование связей во сне & уменьшение контактов измерено; главная роль сна --- гипотеза \\
быстрый сон & работа с отключёнными входами и выходом & режим измерен; назначение неизвестно \\
промывка & расширение межклеточных щелей & противоречивые результаты \\
местный сон & отдельные участки проваливаются в тишину & измерено \\
\bottomrule
\end{tabular}
\caption{Что делают нейроны во сне.}
\label{tab:sleep}
\end{table}

Таблица~\ref{tab:sleep} сводит главу. Сон оказался не паузой в работе машины, а её обслуживанием на ходу: переключение режимов теми же широковещательными каналами, генератор из тех же деталей, что ритм шагов, перемотка дня для медленной памяти и перенормировка весов. Днём машина пишет, ночью переписывает и приводит в порядок записанное.

\section{Задачи}

\begin{exercise}[\lvlA]
\label{ex:20:1}
\emph{Реле без суточного ритма.} Пусть пороги постоянны: $H=0.67$, $L=0.17$ (средние значения порогов модели рис.~\ref{fig:twoprocess}). Выведите из~\eqref{eq:sleeppressure} длительности бодрствования и сна и период генератора при $\tau_r=18.2$~ч, $\tau_d=4.2$~ч. Почему модель с качающимися порогами всё же выходит на $16.1+8.0=24$~ч? К какому прибору главы~\ref{sec:chemoutput} относится этот эффект?
\end{exercise}

\begin{exercise}[\lvlA]
\label{ex:20:2}
\emph{Долг сна.} Сон, начатый при давлении $S_0$, длится $t_s=\tau_d\ln(S_0/L)$, $\tau_d=4.2$~ч, $L$ --- нижний порог при пробуждении. (а)~После $40$~ч без сна $S_0=0.90$, а сон длился $8.8$~ч. Каков был $L$? Сколько длился бы сон при обычном $L\approx0.092$? (б)~За ночь площадь синаптических контактов уменьшается на $18\%$. На сколько должны вырасти веса за день?
\end{exercise}

\begin{exercise}[\lvlB]
\label{ex:20:3}
\emph{Проектирование порогов.} Подберите постоянные пороги $H$ и $L$, при которых реле~\eqref{eq:sleeppressure} с $\tau_r=18.2$~ч и $\tau_d=4.2$~ч даёт ровно $16$~ч бодрствования и $8$~ч сна. Чем такая машина хуже машины с качающимися порогами, если однажды ночью её разбудили через $4$~часа?
\end{exercise}

\begin{exercise}[\lvlB]
\label{ex:20:4}
\emph{Период медленных качелей.} В модели~\eqref{eq:updown} $F(x)=1/\bigl(1+e^{-(x-\theta)/k}\bigr)$, $w=5$, $I=2$, $\theta=2.5$, $k=0.25$, $b=5$, $\tau\ll\tau_a=0.3$~с. (а)~Найдите точки поворота $a_{\text{в}}$ и $a_{\text{н}}$ кривой $r=F(wr+I-a)$, где исчезают работа и тишина. (б)~Приняв $r\approx1$ в работе и $r\approx0$ в тишине, найдите период и долю работы.
\end{exercise}

\begin{exercise}[\lvlB]
\label{ex:20:5}
\emph{Когда качели выключаются.} В модели задачи~\ref{ex:20:4} неподвижная точка лежит на пересечении кривой $a=wr+I-F^{-1}(r)$ с прямой $a=br$. Колебания идут, пока пересечение на средней ветви, между точками поворота. При каких $b$ это так? Как, меняя $b$, \nt{ACET} переключает генератор в усилитель?
\end{exercise}

\begin{exercise}[\lvlB]
\label{ex:20:6}
\emph{Моделирование: долг или фаза.} В \texttt{sleep\_models.py} возьмите первое пробуждение после $48$~ч и держите машину бодрствующей $D=24$, $32$, $40$, $48$, $64$~ч (\texttt{stay\_awake}, \texttt{days=8}). Сколько длится восстановительный сон при каждом $D$? Что определяет его длительность?
\end{exercise}

\begin{exercise}[\lvlB]
\label{ex:20:7}
\emph{Моделирование: забывание и перемешивание.} Линейный нейрон $y=\mathbf w\cdot\mathbf x$, $n=40$, учится по правилу $\mathbf w\leftarrow\mathbf w-0.5\,(y-y^*)\,\mathbf x$. Задачи A и B --- по $10$ образов, $x_j\sim\mathcal N(0,1/n)$, $y^*=\pm1$. Какова среднеквадратичная ошибка на A после $200$ эпох A и затем $200$ эпох B? А после $200$ эпох A и B вперемешку?
\end{exercise}

\begin{exercise}[\lvlC]
\label{ex:20:8}
\emph{Исследование: перемешивание или перенормировка?} Нейрон задачи~\ref{ex:20:7} с порогом; две ночные процедуры: (i)~все веса умножаются на $0.82$; (ii)~старые образы повторяются вперемешку с новыми. Что каждая делает с отношениями весов и с ответами нейрона? Как гипотезы различить по размерам синапсов после сна?
\end{exercise}
